\subsection*{Supplementary Note \thesuppnote: numerical specifications and validation}
\label{si:numerical}


Amplitude counts refer to shared participant--component observations. Acceptance criteria and comparison metrics follow Methods (``Certification contract''): \SI{0.1}{\uV} for amplitudes; identical cluster partitions and \(t\)-maps within \(10^{-12}\) relative, using shared subject-level contrasts and adjacency.

\subsubsection*{Reference implementation and trial identities}

The versioned MMN implementation is \path{tools/validation/validate_mmn_group.py} at tag \texttt{v0.3.2}; Supplementary Table~\ref*{tab:spec} gives the settings for all five components. The code removes \path{HEOG_left}, \path{HEOG_right} and \path{VEOG_lower} before average referencing. It assigns the 30 EEG channels a case-insensitive \path{standard_1020} montage (so FP1/FP2 match Fp1/Fp2), replacing the input file's 33-point digitization. Event samples are rounded from onset times at the analysis sampling rate; samples at or before zero or at or beyond the recording length are discarded. Both measurement-window endpoints are included. The other four components use their own event decoding and \texttt{reject=None}; their minimum retained-trial counts are included in the table.

\begin{table}[!htbp]
\centering\footnotesize
\caption{Numerical specification of the five certified analyses. All: average reference; zero-phase FIR band-pass (\mne{}: lower transition width 0.1\,Hz, upper transition width by \mne{}'s default rule, a quarter of the cutoff and at least 2\,Hz); resampling to 256\,Hz; baseline-corrected epochs; amplitude = mean over the region of interest (ROI) and window of the condition difference wave; group test: two-sided one-sample \(t\)-test of the ROI amplitudes against zero (the stated direction is the expected sign) and a one-sided cluster-permutation test in that direction on the same window (\(B = 5000\), cluster-forming threshold the one-sided \(t_{0.95}\) quantile, shared channel adjacency). Software: \mne{} 1.12.1 (NumPy 2.4.6, SciPy 1.17.1), MNE-BIDS-Pipeline 1.10.1; \eeglab{} development version (git 4207693, 2026-05-28) and \fieldtrip{} 20170830 under GNU Octave 10.3.0.}
\label{tab:spec}
\resizebox{\linewidth}{!}{\begin{tabular}{lccccc}
\toprule
 & MMN & P3b & N170 & ERN & N400 \\
\midrule
participants (shared) & 40 (38) & 20 (20) & 20 (20) & 20 (14) & 20 (20) \\
contrast & deviant \(-\) standard & rare \(-\) frequent & faces \(-\) cars & error \(-\) correct & unrelated \(-\) related \\
time lock & stimulus & stimulus & stimulus & response & stimulus \\
band-pass (Hz) & 0.1--30 & 0.1--30 & 0.1--40 & 0.1--30 & 0.1--30 \\
epoch (s) & \(-0.2\)--0.5 & \(-0.2\)--0.8 & \(-0.2\)--0.5 & \(-0.4\)--0.6 & \(-0.2\)--0.8 \\
baseline (s) & \(-0.2\)--0 & \(-0.2\)--0 & \(-0.2\)--0 & \(-0.4\)--\(-0.2\) & \(-0.2\)--0 \\
rejection & \SI{100}{\uV} peak-to-peak & none & none & none & none \\
minimum retained trials & 50 deviant, 150 standard & 30 rare, 100 frequent & 30 per condition & 15 per condition & 30 per condition \\
ROI & Fz, FCz, Cz & Fz, Cz, Pz, CPz & PO7, PO8, P7, P8 & FCz, Fz, Cz & CPz, Cz, Pz \\
window (s) & 0.10--0.25 & 0.30--0.50 & 0.13--0.20 & 0--0.10 & 0.30--0.50 \\
direction & negative & positive & negative & negative & negative \\
\bottomrule
\end{tabular}}
\end{table}

A post-release check traced trial identities for all 120 available subject-component records: 40 MMN and 20 each for P3b, N170, ERN and N400, including the eight records excluded by the original trial-count criteria. The original AEA subject functions were regenerated without changing their numerical settings; their historical trial identities had not been saved. Reference identities were recovered from retained filtered recordings, pre-rejection epochs and cleaned epochs. Each identity was mapped to the original event-table row using its onset, event code and condition, with annotation labels, selection indices and drop reasons checked against that mapping. Retained memberships and sample indices were then compared, and regenerated counts and ROI amplitudes were checked against the original stored results. This check concerns AEA versus MNE-BIDS-Pipeline; the \eeglab{} and \fieldtrip{} trial comparisons remain count-based.

\subsubsection*{Reference and cross-toolbox comparisons}

All 112 shared participant--component amplitudes pass the reference-pipeline comparison (Supplementary Table~\ref*{tab:cert}). The maximum difference occurred in the MMN participant with the trial-membership mismatch described below. Across all channels of the full difference waveforms, median participant-level root-mean-square differences are 0.011 (MMN), 0.000 (P3b), 0.027 (N170), 0.029 (ERN) and \SI{0.000}{\uV} (N400); per-participant maxima are 0.138, 0.010, 0.099, 0.137 and \SI{0.212}{\uV}, respectively, with isolated single-sample differences reaching \SI{1.6}{\uV}. This comparison tests composition of shared \mne{} kernels, as defined in Methods (``Certification contract''); independent-toolbox comparisons follow below.

\begin{table}[!htbp]
\centering\small
\caption{Per-subject certification of the five reference implementations against MNE-BIDS-Pipeline 1.10.1 in the configuration of Supplementary Table~\ref*{tab:spec}. \(n\): shared subjects; CCC: Lin's concordance correlation coefficient; max \(\absdiff\): largest absolute per-subject difference, computed from the stored per-subject amplitudes (three decimals); recomputed from the retained waveform summaries it is 0.001 (P3b) and 0.007\,\uV{} (N400), both within tolerance.}
\label{tab:cert}
\begin{tabular}{lrrrr}
\toprule
component & \(n\) & CCC & max \(\absdiff\) (\uV) & within \SI{0.1}{\uV} \\
\midrule
MMN & 38 & 0.999747 & 0.090 & 38/38 \\
P3b & 20 & 1.000000 & 0.000 & 20/20 \\
N170 & 20 & 0.999990 & 0.011 & 20/20 \\
ERN & 14 & 0.999993 & 0.032 & 14/14 \\
N400 & 20 & 1.000000 & 0.006 & 20/20 \\
\bottomrule
\end{tabular}
\end{table}
The trial-identity check completed all 120 records and recovered 56,858 candidate event identities. Of the 112 shared records in Supplementary Table~\ref*{tab:cert}, 111 had identical retained memberships: all 74 without rejection and 37 of 38 MMN records. The sole difference was one deviant in MMN sub-030 (event-table row 927, zero-based), whose P9 peak-to-peak amplitude was \SI{101.718}{\uV} in the regenerated AEA epoch and \SI{91.989}{\uV} in the retained reference epoch, straddling the \SI{100}{\uV} rejection threshold. Its event position differed by one sample at 256\,Hz. Across all candidate events, 590 sample positions differed by one sample and the rest agreed. These observations identify the differing trial and timing offsets without isolating their contribution to the amplitude discrepancy. All 112 regenerated AEA records reproduced the original counts and three-decimal ROI amplitudes; the eight original exclusions were preserved.

\paragraph{Toolbox defaults.}

All ten component-by-toolbox analyses reproduced the certified effect direction and significance decision on shared subjects, despite per-subject differences exceeding tolerance (Supplementary Table~\ref*{tab:crosstoolbox}A). The independent \eeglab{} and \fieldtrip{} arms use GNU Octave, two filter implementations and three resampling methods, as detailed in Methods (``Cross-toolbox arms''). Across all 17 runs, five amplitudes changed sign relative to the certified values: P3b sub-014 (\(-0.154\) to \(+0.115\)~\uV) and MMN sub-003 (\(-0.091\) to \(+0.098\)~\uV) at the \fieldtrip{} default; P3b sub-006 (\(-0.760\) to \(+0.042\)~\uV) and sub-014 (\(-0.154\) to \(+0.125\)~\uV) in the width-only \fieldtrip{} configuration; and MMN sub-023 (\(+0.059\) to \(-0.087\)~\uV) under \eeglab{}'s absolute-amplitude criterion.

\begin{table}[!htbp]
\centering
\small
\caption{(A) Per-subject agreement of two independently implemented toolboxes with the certified values at each toolbox's default configuration (before pinning; \eeglab{} default \texttt{pop\_eegfiltnew}; \fieldtrip{} \texttt{firws} at its default transition width, not its general default Butterworth filter; MMN uses the peak-to-peak criterion matched to the reference). $|$grand $\Delta|$ is the absolute difference between the toolbox and certified grand means over the $n$ shared subjects, as a percentage of the certified grand mean; $t$ and $p$ are the toolbox arm's two-sided one-sample test of the $n$ shared subjects' amplitudes against zero ($n-1$ degrees of freedom); same conclusion means that the effect direction and the $p<0.05$ decision equal the reference's. (B) Controlled filter comparisons in \fieldtrip{}, reporting maximum absolute per-subject ROI amplitude differences in \uV{}.}
\label{tab:crosstoolbox}
\textbf{A. Default toolbox configurations}\par\smallskip
\resizebox{\linewidth}{!}{\begin{tabular}{llrrrrrrrc}
\toprule
toolbox & component & $n$ & $|$grand $\Delta|$ (\%) & max $|\Delta|$ (µV) & CCC & $\leq$0.1\,µV & $t$ & $p$ (two-sided) & same conclusion \\
\midrule
EEGLAB & MMN & 38 & 1.844 & 0.233 & 0.9953 & 33/38 & -7.68 & $3.5\times10^{-9}$ & yes \\
EEGLAB & P3b & 20 & 0.520 & 0.050 & 0.9999 & 20/20 & 4.14 & $5.6\times10^{-4}$ & yes \\
EEGLAB & N170 & 20 & 3.320 & 0.161 & 0.9979 & 16/20 & -4.63 & $1.8\times10^{-4}$ & yes \\
EEGLAB & ERN & 14 & 0.338 & 0.210 & 0.9997 & 12/14 & -6.31 & $2.7\times10^{-5}$ & yes \\
EEGLAB & N400 & 20 & 0.512 & 0.069 & 0.9999 & 20/20 & -6.38 & $4.1\times10^{-6}$ & yes \\
FieldTrip & MMN & 38 & 1.177 & 0.189 & 0.9939 & 30/38 & -7.75 & $2.9\times10^{-9}$ & yes \\
FieldTrip & P3b & 20 & 3.779 & 0.721 & 0.9928 & 10/20 & 4.43 & $2.9\times10^{-4}$ & yes \\
FieldTrip & N170 & 20 & 1.385 & 0.117 & 0.9990 & 18/20 & -4.39 & $3.1\times10^{-4}$ & yes \\
FieldTrip & ERN & 14 & 2.275 & 0.366 & 0.9979 & 3/14 & -6.29 & $2.8\times10^{-5}$ & yes \\
FieldTrip & N400 & 20 & 0.823 & 0.229 & 0.9982 & 10/20 & -6.46 & $3.5\times10^{-6}$ & yes \\
\bottomrule
\end{tabular}}

\par\medskip
\textbf{B. FieldTrip filter settings}\par\smallskip
\begin{tabular}{lrrr}
\toprule
Component & Default & Width only & Cutoff + width matched \\
\midrule
P3b & 0.7214 & 0.8023 & 0.0496 \\
N400 & 0.2287 & n.m. & 0.0500 \\
ERN & 0.3658 & n.m. & 0.0697 \\
N170 & 0.1170 & n.m. & 0.0993 \\
MMN & 0.1888 & n.m. & 0.1838 \\
\bottomrule
\end{tabular}
\par\smallskip
\begin{minipage}{\linewidth}\footnotesize
Width only fixes the transition width to 0.1\,Hz; cutoff + width also interprets the stated cutoff as the passband edge. n.m. marks an unmeasured configuration; only P3b has an intermediate measurement. The tolerance is \SI{0.1}{\uV}. All four components without trial rejection meet it after matching; MMN remains above tolerance and its residual is unresolved. The default-to-matched comparison changes two parameters; the cutoff-only attribution rests on the P3b intermediate comparison. Values are from the result records listed in Supplementary Table~\ref*{tab:files}.
\end{minipage}
\end{table}
\paragraph{Differences relative to between-subject variance.}

The paired-difference standard deviations are 1.3--11.3\,\% of the between-subject standard deviations. Supplementary Table~\ref*{tab:ba} gives each comparison's bias, dispersion, Bland--Altman limits and illustrative sample-size inflations under the two measurement-error scenarios defined in Methods (``Size of toolbox-induced differences'').

\begin{table}[!htbp]
\centering\small
\caption{Default-toolbox agreement with certified values. Bias (toolbox minus certified), paired-difference standard deviation \(s_d\), reference between-subject standard deviation \(s_b\), and the two sample-size inflation scenarios follow Methods (``Size of toolbox-induced differences''). Limits of agreement are bias \(\pm 1.96\,s_d\). Source: \texttt{tools/benchmark/reports/attenuation.json}.}
\label{tab:ba}
\resizebox{\linewidth}{!}{\begin{tabular}{llrrrlrrr}
\toprule
toolbox & component & \(n\) & bias (\uV) & \(s_d\) (\uV) & limits (\uV) & \(s_b\) (\uV) & equal split (\%) & one-sided (\%) \\
\midrule
EEGLAB & ERN & 14 & +0.018 & 0.080 & -0.138 to +0.175 & 3.23 & 0.03 & 0.06 \\
EEGLAB & MMN & 38 & -0.015 & 0.063 & -0.139 to +0.108 & 0.66 & 0.46 & 0.92 \\
EEGLAB & N170 & 20 & -0.039 & 0.065 & -0.167 to +0.089 & 1.19 & 0.15 & 0.30 \\
EEGLAB & N400 & 20 & -0.015 & 0.027 & -0.069 to +0.039 & 2.09 & 0.01 & 0.02 \\
EEGLAB & P3b & 20 & -0.009 & 0.027 & -0.061 to +0.043 & 1.81 & 0.01 & 0.02 \\
FieldTrip & ERN & 14 & +0.124 & 0.166 & -0.201 to +0.449 & 3.23 & 0.13 & 0.26 \\
FieldTrip & MMN & 38 & -0.010 & 0.073 & -0.152 to +0.133 & 0.66 & 0.61 & 1.22 \\
FieldTrip & N170 & 20 & +0.016 & 0.050 & -0.081 to +0.114 & 1.19 & 0.09 & 0.18 \\
FieldTrip & N400 & 20 & -0.024 & 0.122 & -0.264 to +0.215 & 2.09 & 0.17 & 0.34 \\
FieldTrip & P3b & 20 & +0.064 & 0.205 & -0.337 to +0.465 & 1.81 & 0.64 & 1.27 \\
\bottomrule
\end{tabular}}
\end{table}
\subsubsection*{Numerical conventions and implementation checks}

Supplementary Tables~\ref*{tab:crosstoolbox}B and~\ref*{tab:cluster}B report the controlled comparisons.

\paragraph{Rejection criterion.} MMN is the only certified component with trial rejection. Replacing \eeglab{}'s absolute-amplitude reading of ``\SI{\pm 100}{\uV}'' with \mne{}'s peak-to-peak reading reduced the mean absolute difference from certified values 2.24-fold (Supplementary Table~\ref*{tab:instances}) and increased the within-tolerance count from 22 to 33 of 38. Even with the matched criterion, \eeglab{} retained 130 deviant trials for the first participant versus \mne{}'s 117. One candidate cause, not separately isolated, is that \mne{}'s automatic upper transition band admits more high-frequency content and therefore rejects more trials.

\paragraph{Filter specification.} Pinning transition width alone increased the \fieldtrip{} P3b discrepancy; changing the cutoff convention at that fixed width removed most of it (Supplementary Table~\ref*{tab:crosstoolbox}B). The latter comparison yielded Pearson and Spearman correlations of 1.0000 and all 20 observations within tolerance. Reading the stated cutoff as the passband edge (\mne{} and \eeglab{}) or the \SI{-6}{dB} point (\fieldtrip{}) places the high-pass \SI{-6}{dB} point at 0.05 or 0.1\,Hz, a factor of two apart. Thus ``0.1--30\,Hz zero-phase FIR'' specifies different filters. With the specification pinned for all five components, the four components without rejection gave 74 of 74 comparisons within tolerance (20 distinct subjects), versus 41 of 74 at \fieldtrip{} defaults. \eeglab{} gave 68 of 74 at its defaults and was not pinned further. MMN improved from 30 to 34 of 38 within tolerance but remained uncertified against the independent toolbox. Its residual is consistent with rejection coupling and the unpinned \mne{} upper transition width; neither was isolated for MMN.

\paragraph{Cluster-forming threshold.} In the full-epoch N400 comparison, the historical threshold error changed the leading summed \(t\) by 6.3\,\%. Supplementary Table~\ref*{tab:cluster}B gives the cluster counts and extents under the erroneous, specified and matched thresholds, including the unchanged leading-cluster significance.

\paragraph{Observed cluster statistic.}

At the shipped configuration (Supplementary Table~\ref*{tab:cluster}A), both implementations gave identical partitions and leading summed \(t\) statistics to every printed digit. The observed \(t\)-maps differed by at most \(1.78\times10^{-15}\), with the implementations summing in different orders. \mne{} reproduced the committed cluster counts and \(p\)-values exactly. \fieldtrip{}'s leading-cluster \(p\)-values differed by at most one permutation count (0.0002 at \(B = 5000\)); the one non-leading cluster (P3b) also differed by 0.0002, within Monte-Carlo error. Its permutations are unseeded, so its \(p\)-values are excluded from exactness statements. Matched full-epoch comparisons also agreed on observed statistics and partitions for all four components. Leading-cluster \(p\)-values were 0.0002--0.0004 (the Monte-Carlo floor or one count above), except for N170: 0.0124 in \mne{} and 0.0128 in \fieldtrip{}, two counts apart and a quarter of one estimate's standard error.

\begin{table}[!htbp]
\centering\small
\caption{(A) Cluster-permutation certification at the shipped configuration (component window, one-sided \(t_{0.95}\) threshold, \(B = 5000\), shared 30-channel adjacency). Leading summed \(t\) and extent are identical in the two implementations; \(t\)-map max \(\absdiff\) is the largest element-wise difference; partitions were compared as sets of member samples. \(p\): \mne{} / \fieldtrip{}. (B) Separate full-epoch N400 threshold-sensitivity comparison; counts are suprathreshold clusters and extent is in channel--time samples.}
\label{tab:cluster}
\textbf{A. Shipped component-window configuration}\par\smallskip
\begin{tabular}{lrrrrrll}
\toprule
component & window (ms) & threshold & clusters & leading \(\sum t\) & extent & \(t\)-map max \(\absdiff\) & \(p\) \\
\midrule
N400 & 301--500 & \(-1.729\) & 1 / 1 & \(-2102.71\) & 518 & \(1.8\times10^{-15}\) & 0.0002 / 0.0002 \\
ERN & 0--98 & \(-1.771\) & 1 / 1 & \(-1014.83\) & 257 & \(1.8\times10^{-15}\) & 0.0004 / 0.0006 \\
N170 & 133--199 & \(-1.729\) & 1 / 1 & \(-645.14\) & 166 & \(1.8\times10^{-15}\) & 0.0002 / 0.0004 \\
P3b & 301--500 & \(+1.729\) & 2 / 2 & \(+1368.89\) & 415 & \(1.8\times10^{-15}\) & 0.0002 / 0.0004 \\
\bottomrule
\end{tabular}

\par\medskip
\textbf{B. Full-epoch N400 threshold comparison}\par\smallskip
\begin{tabular}{>{\raggedright\arraybackslash}p{6.2cm}rrr}
\toprule
Implementation and configuration & Quantile & \shortstack{Suprathreshold\\clusters} & \shortstack{Leading-cluster\\extent} \\
\midrule
\mne{}, historical error & 0.975 & 14 & 1268 \\
\fieldtrip{}, specified one-sided threshold & 0.95 & 18 & 1419 \\
\fieldtrip{}, matched historical error & 0.975 & 14 & 1268 \\
\bottomrule
\end{tabular}
\par\smallskip
\begin{minipage}{\linewidth}\footnotesize
The leading cluster has the maximum absolute cluster statistic, not necessarily the largest extent. All three leading clusters in B remain significant; the counts are not counts of significant clusters. The historical 0.975 quantile was a configuration error for the intended one-sided test. Matching that error across toolboxes establishes agreement, not correctness. The released configuration uses 0.95; its component-window certification in A is distinct from the full-epoch comparison in B. Source records are listed in Supplementary Table~\ref*{tab:files}.
\end{minipage}
\end{table}
\paragraph{Specification and implementation issues.}

Supplementary Table~\ref*{tab:instances} records all eight observed issues: three controlled effects of unstated conventions, one parameter whose pinning did not reduce disagreement, one implementation error despite an adequate specification, and three associated or unisolated observations. The rejection and cutoff comparisons changed one convention at a time; the Welch comparison changed two together.

\begin{table}[!htbp]
\centering\footnotesize
\caption{Eight observed specification or implementation issues. Consequences are per-subject amplitude differences unless stated.}
\label{tab:instances}
\begin{tabular}{>{\raggedright\arraybackslash}p{2.5cm}>{\raggedright\arraybackslash}p{4.0cm}>{\raggedright\arraybackslash}p{4.8cm}>{\raggedright\arraybackslash}p{2.5cm}}
\toprule
phrase & the two readings & measured consequence & type \\
\midrule
``\SI{\pm 100}{\uV} rejection'' & absolute amplitude vs peak-to-peak & mean \(\absdiff\) 0.0886 vs \SI{0.0396}{\uV} (MMN, \eeglab{}) & controlled effect \\
``0.1--30\,Hz FIR'' (cutoff) & stated frequency at the \SI{-6}{dB} point vs at the passband edge & P3b max \(\absdiff\) 0.8023 to \SI{0.0496}{\uV} at fixed width & controlled effect \\
transition width & 0.2 vs 0.1\,Hz at both edges & P3b 0.7214 to \SI{0.8023}{\uV}: pinning did not reduce the difference & no reduction \\
``\(\alpha = 0.05\), one-sided'' & \(t_{0.975}\) used vs \(t_{0.95}\) specified & 14 vs 18 suprathreshold clusters, leading extent 1268 vs 1419 (N400, full epoch); leading cluster significant in both & implementation error (specification adequate) \\
``Welch, Hamming window'' & symmetric window without demeaning vs periodic window with per-segment demeaning & band power 0.44\,\% to \(6.9\times10^{-16}\) relative (joint effect of both conventions) & controlled effect \\
reference scheme & mastoid vs average reference & 7 of 9 unconstrained generations chose mastoid: grand-mean error 0.894 vs \SI{0.123}{\uV} (different cohorts; Supplementary Note~\ref*{si:generation}) & associated observation \\
RANSAC / automated rejection & on vs off & retained trial sets differ; not isolated numerically & unisolated risk \\
window-edge alignment & sample at the edge included vs excluded & two readings observed in generated programs (Supplementary Note~\ref*{si:generation}); no isolated consequence & unisolated risk \\
\bottomrule
\end{tabular}
\end{table}
\subsubsection*{Spectral and known-effect checks}

\paragraph{Spectral implementation.}

\mne{} Welch band power on 12 raw eyes-closed recordings agreed with a from-scratch NumPy implementation across four bands after matching the periodic window and per-segment demeaning: maximum relative error \(6.9\times10^{-16}\) for Hamming and \(5.2\times10^{-16}\) for Hann, the released skill and recipe default. Thus the comparison covers both the original configuration and the released default (Methods, ``Spectral, figure-fidelity and known-effect checks''). The original scratch configuration (symmetric window, no demeaning), retained separately, had maximum relative errors of 0.44, 0.37, 0.25 and 0.14\,\% for delta, theta, alpha and beta. Changing both conventions together establishes their joint effect. Replacing Welch with multitaper changed band power by 0.7--4.6\,\% for Hamming and 0.8--5.8\,\% for Hann (medians across recordings, per band), demonstrating estimator sensitivity.

\paragraph{Known effects on public data.}

The unchanged project-authored resting-state script, rerun after release on all 109 PhysioNet participants, showed differences under the retained ICA configuration. Posterior alpha power was higher eyes-closed than eyes-open in 106 of 109 participants (\(t(108) = 5.71\), \(p = 1.0\times10^{-7}\)); Lempel--Ziv complexity and permutation entropy (\(p = 2.7\times10^{-27}\)) were higher eyes-open. With four microstate classes imposed, mean total global explained variance was 0.665 and pooled median class duration 104.1\,ms. For S001--S020, the rerun reproduced released group statistics (alpha \(t(19) = 3.61\), \(p = 0.0019\)). Of 260 stored participant measurements, 257 agreed within \(3\times10^{-14}\) relative; one eyes-open Lempel--Ziv value differed by \(1.4\times10^{-5}\), and two microstate durations by 0.06 and 0.05\,ms.

Omitting ICA in the same 109 participants reduced the eyes-open-minus-eyes-closed Lempel--Ziv difference from 0.0309 to 0.0049. The corresponding tests were \(t(108) = 7.24\), \(p = 6.9\times10^{-11}\) with ICA and \(t(108) = 0.96\), \(p = 0.34\) without ICA. Without ICA, permutation entropy remained higher eyes-open (difference 0.0268; \(t(108) = 11.98\), \(p = 1.5\times10^{-21}\)). Posterior alpha remained higher eyes-closed in 106 of 109 participants (\(t(108) = 5.69\), \(p = 1.1\times10^{-7}\)). The Lempel--Ziv state contrast was therefore sensitive to ICA in these data. This comparison does not distinguish removal of physiological artifacts from differences introduced by ICA processing.

The original adapted N170 run-01 analysis on 16 ds000117 participants gave faces-minus-scrambled amplitude \SI{-0.641}{\uV} (SEM 0.251; \(t(15) = -2.56\), \(p = 0.022\), \(d_z = -0.64\)) and one significant cluster in the recorded run (\(p = 0.029\)). Requiring the recipe's 30 trials per condition retained 15 participants and the ROI effect (mean -0.70\,\uV{}, \(t(14) = -2.70\), \(p = 0.017\)); the subset's cluster test was not repeated.

Adding only the 34.5\,ms trigger shift to the original route gave a mean of \SI{-1.472}{\uV} (SEM 0.273; \(t(15) = -5.38\), \(p = 7.6\times10^{-5}\), \(d_z = -1.35\)), negative in 15 of 16 participants. Its minimum cluster \(p\)-value was 0.0008. The six-run route with the same shift, named-channel exclusions and fixed ROI gave \SI{-1.427}{\uV} (SEM 0.117; \(t(15) = -12.21\), \(p = 3.4\times10^{-9}\), \(d_z = -3.05\)). All 16 amplitudes were negative, and the minimum cluster \(p\)-value was 0.0006. Using participant-specific ROIs with the same six-run data gave \SI{-1.534}{\uV} (SEM 0.124; \(t(15) = -12.41\), \(p = 2.7\times10^{-9}\), \(d_z = -3.10\)); no additional cluster test was recorded for that ROI summary. The original route reproduced the released three-decimal participant amplitudes exactly. The delay-only comparison isolates the trigger-shift change; the six-run comparison also changes run coverage, channel exclusion and ROI definition. These adapted checks do not certify the recipes as written.

\subsubsection*{Evidence coverage and limits}

Supplementary Table~\ref*{tab:knowninput} reports the post-release known-input checks: eight of nine interface-tested skills agreed with analytic or explicitly coded references. The released QC rule failed on artifact-free 10\,Hz sinusoids. The revised rule passed the analytic checks and located injected faults, including those in three MNE sample channels. Supplementary Table~\ref*{tab:skill-coverage} records the unchanged release inventory and its scope.

\begin{table}[!htbp]
\centering\footnotesize
\caption{Known-input checks dated 3 October 2026, outside the \texttt{v0.3.2} snapshot, for skills classified L0 in Supplementary Table~\ref*{tab:skill-coverage}. Unless stated, the result is computed for each checked output as the largest absolute difference divided by the largest absolute reference value, and the table gives the largest such ratio across outputs; a control with an all-zero reference is compared by absolute difference (0). Records: \texttt{tools/validation/library\_numerics\_2026-10-03.json} (with the initial run \texttt{library\_numerics\_initial\_2026-10-03.json}), \texttt{qc\_revision\_2026-10-03.json} and \texttt{qc\_sample\_revision\_2026-10-03.json} at commit \texttt{546db87}; runner \texttt{tools/validation/library\_numerics.py}.}
\label{tab:knowninput}
\begin{tabular}{>{\raggedright\arraybackslash}p{2.3cm}>{\raggedright\arraybackslash}p{4.6cm}>{\raggedright\arraybackslash}p{4.0cm}>{\raggedright\arraybackslash}p{3.2cm}}
\toprule
skill & operation checked & reference & result \\
\midrule
eeg-tfr & zero-mean Morlet power, inter-trial coherence, baseline units & explicit Gaussian wavelet, direct convolution & \(2.4\times10^{-15}\) \\
eeg-connectivity & debiased squared weighted phase-lag index, three edges, 8--13\,Hz & direct DFT and moment formula & \(6.9\times10^{-15}\) \\
eeg-ica & extended Infomax on a known three-source mixture; component removal & known sources and clean mixture & minimum \(|r|\) 0.9999; cleaned relative RMSE 0.021 \\
eeg-source & MNE and dSPM, fixed orientation, spherical model & direct regularized linear inverse & \(2.9\times10^{-15}\) \\
eeg-decoding & sliding estimator, group-disjoint cross-validation & explicit fold and time loop, same estimator & 0 \\
eeg-complexity & permutation entropy, order 3 & explicit ordinal-pattern counts & \(2.2\times10^{-16}\) \\
eeg-microstate & state-transition probabilities & row-normalized literal adjacent-state counts & 0 \\
eeg-bids & BrainVision round trip of signals, channels and events & original values and metadata & \(3.7\times10^{-8}\) \\
eeg-qc, released rule & flat-signal and peak annotation & artifact-free sinusoids; injected faults & failed: 2 of 3 clean channels flagged, 100 annotations \\
eeg-qc, revised rule & constant spans, rapid changes, low-amplitude warnings & 27 analytic waveforms; injected faults in synthetic and MNE sample data & no clean-waveform candidates; injected faults located \\
\bottomrule
\end{tabular}
\end{table}
\begin{table}[!htbp]
\centering\footnotesize
\caption{Evidence categories assigned to all 22 skills in the v0.3.2 certification inventory, for the operations and configurations tested. Counts refer to skills, not independent tests.}
\label{tab:skill-coverage}
\begin{tabular}{>{\raggedright\arraybackslash}p{2.8cm}>{\raggedright\arraybackslash}p{5.5cm}>{\raggedright\arraybackslash}p{6.3cm}}
\toprule
Evidence category & Skills in the release inventory & What the evidence establishes \\
\midrule
L3: cross-toolbox comparison (4) & eeg-preprocess, eeg-epoch, eeg-erp, eeg-stats & An independently implemented toolbox reproduces the tested per-subject values or observed group statistic within the stated acceptance criterion. \\
\addlinespace[5pt]
L2: cross-implementation comparison (1) & eeg-spectral & A second implementation of the same numerics reproduces the tested output. \\
\addlinespace[5pt]
L1: behavioural check (3) & eeg-audit, eeg-methods-text, eeg-figure & Behaviour is evaluated through seeded-defect detection or checklist coverage. \\
\addlinespace[5pt]
L0: interface test (14) & eeg-recipe, eeg-ica, eeg-tfr, eeg-connectivity, eeg-source, eeg-microstate, eeg-decoding, eeg-complexity, eeg-behavior, eeg-bids, eeg-qc, eeg-group-compare, eeg-report, eeg-pipeline & Automated tests exercise the documented interface, from import and presence checks to synthetic-data runs returning correctly shaped results. These checks do not establish numerical correctness. \\
\bottomrule
\end{tabular}
\par\medskip
\begin{minipage}{\linewidth}\footnotesize
Each skill has one assigned evidence category. The categories do not score overall capability or imply that every lower category was separately tested. Preprocessing, epoching and ERP measurement have end-to-end evidence through the tested ROI amplitudes, not stage-by-stage certification. Independent amplitude certification covers the four components without trial rejection; MMN is not independently certified. The statistics evidence covers the evaluated observed cluster statistics and memberships, not exact Monte-Carlo \(p\)-value agreement. Post-release known-input checks and the revised QC rule (Supplementary Table~\ref*{tab:knowninput}) do not change this release inventory. Source files for the category assignments are listed in Supplementary Table~\ref*{tab:files}.
\end{minipage}
\end{table}


\FloatBarrier
\refstepcounter{suppnote}
